\section*{参考书目}
\phantomsection
\addcontentsline{toc}{section}{参考书目}

\begin{enumerate}
\def\labelenumi{\arabic{enumi}.}
\item
  O. 诺伊格鲍尔，《古代数学史讲义》，第一卷，《前希腊数学》，柏林（施普林格出版社），1934年。
\item
  亚里士多德著作集，由J. A. 史密斯与W. D. 罗斯主编翻译（12卷，牛津，1908—1952年）。
\end{enumerate}

2（bis）。T. L. 希思，《亚里士多德著作中的数学》，牛津（克拉伦登出版社），1949年。

\begin{enumerate}
\def\labelenumi{\arabic{enumi}.}
\setcounter{enumi}{2}
\item
  T. L. 希思，《欧几里得几何原本》十三卷\ldots，3卷本，剑桥，1908年。
\item
  阿基米德全集，共3卷，J. L. 海伯格编，第二版，1913—1915年。
\end{enumerate}

4（bis）。T. L. 希思，《阿基米德方法》，剑桥，1912年。

\begin{enumerate}
\def\labelenumi{\arabic{enumi}.}
\setcounter{enumi}{4}
\item
  J. M. 波亨斯基，《古代形式逻辑》，《逻辑学研究》，阿姆斯特丹（北荷兰出版公司），1951年。
\item
  P. 伯纳，《中世纪逻辑学：其从1250年至约1400年发展概要》，芝加哥，1952年。
\item
  D. 弗朗西斯奇·毛罗利奇，墨西拿修道院院长，著名数学家，算术两卷，威尼斯，1575年。
\item
  伽利略·伽利莱，《著作集》，国家版重印本，20卷，佛罗伦萨（巴巴拉出版社），1929—1939年。
\end{enumerate}

(*) I. NOVAK-GAL，“一致系统模型的构造”，Fund. Math.，37（1950），第87-110页。

\begin{enumerate}
\def\labelenumi{\arabic{enumi}.}
\setcounter{enumi}{8}
\item
  A. 吉拉尔，《代数新发明》，1629年（新版，比伦斯·德·汉，1884年）。
\item
  R. 笛卡尔，《著作集》，C. 亚当与P. 坦纳里编，11卷，巴黎（L. 塞尔夫出版社），1897—1909年。
\item
  B. 帕斯卡，《著作集》，布伦什维格编，14卷，巴黎（阿歇特出版社），1904—1914年。
\item
  G. W. 莱布尼茨：(a)《数学著作》，C. I. 格哈特编，7卷，柏林—哈勒（阿舍尔—施密特），1849—1863年；(b)《哲学著作》，C. I. 格哈特编，7卷，柏林，1840—1890年；(c)《未刊小篇与残篇》，L. 库蒂拉编，巴黎（阿尔康），1903年。
\end{enumerate}

12（bis）。L. 库蒂拉，《根据未刊文献研究莱布尼茨的逻辑》，巴黎（阿尔康出版社），1901年。

\begin{enumerate}
\def\labelenumi{\arabic{enumi}.}
\setcounter{enumi}{12}
\item
  达朗贝尔，《百科全书》，巴黎，1751-1765年，“负数”、“虚数”、“定义”条目。
\item
  C. F. 高斯，《全集》，12卷，哥廷根，1870—1927年。
\item
  N. 罗巴切夫斯基，《泛几何学》，奥斯特瓦尔德经典丛书，第130号，莱比锡（恩格尔曼出版社），1902年。
\item
  G. 布尔：(a)《逻辑的数学分析》，剑桥—伦敦，1847年（=《逻辑著作集》，P. 茹尔丹编，芝加哥—伦敦，1916年，第I卷）；(b)《思维规律的研究》，剑桥—伦敦，1854年（=《逻辑著作集》，第II卷）。
\item
  H. GRASSMANN，全集，6卷，莱比锡（托伊布纳出版社），1894年。
\item
  B. 博尔扎诺，《无穷的悖论》，莱比锡，1851年。
\item
  B. 黎曼，数学著作全集，第二版，莱比锡（托伊布纳出版社），1892年。
\item
  H. HANKEL，《复数系理论》，莱比锡（福斯出版社），1867年。
\item
  A. 德·摩根：(a)《论三段论（III）》，剑桥哲学学会汇刊，第10卷（1858年），第173-230页；(b)《论三段论（IV）及关系逻辑》，剑桥哲学学会汇刊，第10卷（1860年），第331-358页。
\item
  C. S. 皮尔士，(a)《论数学的逻辑》，美国艺术与科学学院院刊，7（1865-1868），第402-412页；(b)《论逻辑代数》，美国数学杂志，3（1880），第49-57页；(c)《论逻辑代数》，美国数学杂志，7（1884），第190-202页。
\item
  E. 施罗德，《逻辑代数讲义》，3卷，莱比锡（托伊布纳出版社），1890年。
\item
  G. 弗雷格，(a)《概念文字——一种模仿算术的纯思维公式语言》，哈雷，1879年；(b)《算术基础》，第2版，附J. L. 奥斯汀的英译本，纽约，1950年；(c)《算术基本定律》，按概念文字推导，2卷，耶拿，1893-1903年。
\item
  G. 康托尔，《论文集》，柏林（施普林格出版社），1932年。
\end{enumerate}

25（续）。G. 康托尔，R. 戴德金，《通信集》，由J. 卡瓦耶斯和E. 诺特编辑，科学与工业丛书，第518号，巴黎（赫尔曼出版社），1937年。

\begin{enumerate}
\def\labelenumi{\arabic{enumi}.}
\setcounter{enumi}{25}
\item
  R. 戴德金，《数学全集》，3卷，不伦瑞克（菲韦格出版社），1932年。
\item
  C. 埃尔米特，T. 斯蒂尔杰斯，《通信集》，2卷，巴黎（戈蒂埃-维拉尔出版社），1905年。
\item
  M. PASCH 和 M. DEHN，《新几何学讲义》，第2版，柏林，（Springer），1926年。
\item
  G. 皮亚诺，(a)《算术原理，新方法阐述》，都灵，1889年；(b)《几何原理，逻辑阐述》，都灵，1889年；(c)《数学公式汇编》，5卷，都灵，1895—1905年。
\item
  D. 希尔伯特，《几何基础》，第7版，莱比锡—柏林（托伊布纳出版社），1930年。
\item
  D. 希尔伯特，《数学论文集》，第 III 卷，柏林（施普林格出版社），1935 年。
\item
  D. 希尔伯特，W. 阿克曼，《理论逻辑基础》，第3版，柏林（施普林格出版社），1949年。
\item
  H. 庞加莱：（a）《科学与假设》，巴黎（弗拉马里翁出版社），1906年；（b）《科学的价值》，巴黎（弗拉马里翁出版社），1905年；（c）《科学与方法》，巴黎（弗拉马里翁出版社），1920年。
\item
  J. 理查德，《数学原理与集合问题》，《纯粹与应用科学总评》，第16卷（1905年），第541-543页。
\item
  R. BAIRE, E. BOREL, J. HADAMARD, H. LEBESGUE，《关于集合论的五封信》，法国数学会通报，33卷（1905年），第261-273页。
\item
  E. 策梅洛，(a)《证明每个集合都可以良序化》，Math. Ann.，59 (1904)，第514-516页；(b)《关于良序可能性的一种新证明》，Math. Ann.，65 (1908)，第107-128页；(c)《关于集合论基础的研究》，Math. Ann.，65 (1908)，第261-281页。
\item
  B. 罗素与A. N. 怀特海，《数学原理》，三卷本，剑桥，1910—1913年。
\item
  L. E. J. 布劳威尔，(a)《直觉主义与形式主义》，美国数学会通报，20卷(1913年)，第81-96页；(b)《论直觉主义数学的基础》，数学年鉴，93卷(1925年)，第244-257页；95卷(1926年)，第453-473页；96卷(1926年)，第451-458页。
\item
  T. 斯科伦，《关于集合论公理化基础的一些评注》，科学报告，第五届斯堪的纳维亚数学家大会，赫尔辛福斯，1922年。
\item
  K. 库拉托夫斯基，《一种从数学推理中消除超限数的方法》，基础数学，5卷(1922年)，第76-108页。
\item
  A. 弗伦克尔：(a)《关于康托-策梅洛集合论的基础》，数学年鉴，86卷(1922年)，第230-237页；(b)《关于集合论基础的十次讲座》，科学与假说丛书第31卷，莱比锡-柏林，1927年；(c)《集合论导论》，第三版，柏林(施普林格出版社)，1928年。
\item
  J. 埃尔布朗，《关于证明论的研究》，华沙科学与文学学会著作集，第二类(1930年)，第33-160页。
\item
  J. 冯·诺依曼：(a)《集合论的公理化》，克雷勒杂志，154卷(1925年)，第219-240页；(b)《集合论的公理化》，数学杂志，27卷(1928年)，第669-752页；(c)《论希尔伯特证明论》，数学杂志，26卷(1927年)，第1-46页。
\item
  K. 哥德尔：(a)《关于《数学原理》及相关系统中的形式不可判定命题》，数学与物理月刊，38 (1931)，第173-198页；(b)《选择公理与广义连续统假设的一致性》，数学研究年鉴，第3号，普林斯顿，1940年。
\item
  M. 佐恩，《关于超限代数中方法的一点评论》，美国数学会通报，41 (1935)，第667-670页。
\item
  A. 海廷，《数学基础研究。直觉主义。证明论》，数学成果，第3卷，柏林（施普林格），1934年。
\item
  G. 根岑，《数学基础研究的当前状况。纯数论无矛盾性证明的新表述》，逻辑研究\ldots，第4辑，莱比锡（希尔策尔），1938年。
\item
  S. 克林，《元数学导论》，纽约，1952年。
\item
  P. J. 科恩，《连续统假设的独立性》，美国国家科学院院刊，50 (1963)，第1143-1148页及51 (1964)，第105-110页。
\end{enumerate}
